\documentclass[12pt]{article}
\usepackage[utf8]{inputenc}
\usepackage[T1]{fontenc}
\usepackage{hyperref}
\title{Phrygian Language Resources for K-12 Curriculum Development}
\author{Compiled by Claude Sonnet 5.5 (claude-sonnet-5-5) for Andr\'as Kornai}
\date{}
\begin{document}
\maketitle

\section{Introduction}
Phrygian (ISO 639-3 xpg, Glottocode phry1239) is an extinct Indo-European language of Anatolia. It is known from Old Phrygian inscriptions (about 395 in the source found) and from Neo-Phrygian funerary texts in the Greek alphabet. No school materials or revitalization programme turned up in the searches. The resources below are scholarly.

\section{Primary Educational Resources}
None found. Any classroom use would have to be adapted from the scholarly sources below.

\section{Linguistic and Reference Materials}
\begin{itemize}
\item A Phrygian overview by Brent Vine and colleagues' circle is not claimed here. The course-hosted introductory PDF found is at \url{https://www.uvm.edu/~jbailly/courses/Indo-European/EggyPhrygian%20(1).pdf}.
\item Calder's \emph{Corpus Inscriptionum Neo-Phrygiarum} II and III in the Journal of Hellenic Studies: \url{https://resolve.cambridge.org/core/journals/journal-of-hellenic-studies/article/corpus-inscriptionum-neophrygiarumii/81BE82C4646F67EC61EB5AC1F65632F7} and \url{https://resolve.cambridge.org/core/journals/journal-of-hellenic-studies/article/corpus-inscriptionum-neophrygiarum-iii/BABC2A2F5103F5057F70DAE520933142}.
\item Further Neo-Phrygian inscriptions: \url{https://resolve.cambridge.org/core/journals/journal-of-hellenic-studies/article/abs/new-neophrygian-inscriptions/D4413A3A4728D10F0A3DAA8425BF4224}.
\item Obrador Cursach's dissertation, including a lexicon and catalogue: \url{https://tesisenred.net/handle/10803/650834?show=full}.
\item De Hoz, Greek--Phrygian contact in the Neo-Phrygian corpus: \url{https://produccioncientifica.usal.es/documentos/659fe270aaaaaa46358cc728}.
\item A paper on Phrygian from the New Bulgarian University repository: \url{https://eprints.nbu.bg/id/eprint/2641/1/8.pdf}.
\item Overview: \url{https://wikipedia.com/wiki/ISO_639:xpg}.
\end{itemize}

\section{Religious and Community Resources}
Nothing found. The Neo-Phrygian texts are mostly funerary curse formulas.

\section{Digital and Audio Resources}
No audio or software was found. The search snippets did not confirm an online corpus beyond the sources above.

\section{Curriculum Development Context}
There are no native speakers and no known revival effort. The corpus is short, so a curriculum could only cover a handful of words and formulas, within a broader unit on ancient Anatolia and Indo-European languages.

\end{document}